\documentclass{article}

\usepackage{amsmath,amssymb}
\usepackage{xurl}
\usepackage[hidelinks]{hyperref}
\setlength{\emergencystretch}{2em}

% Shared commands for notes in docs/paper-gaps/.

\DeclareUnicodeCharacter{2080}{\ifmmode{}_{0}\else\textsubscript{0}\fi}
\DeclareUnicodeCharacter{2081}{\ifmmode{}_{1}\else\textsubscript{1}\fi}
\DeclareUnicodeCharacter{2082}{\ifmmode{}_{2}\else\textsubscript{2}\fi}
\DeclareUnicodeCharacter{2099}{\ifmmode{}_{n}\else\textsubscript{n}\fi}

\newcommand{\Lean}{\textsc{Lean}}
\ExplSyntaxOn
\NewDocumentCommand{\leanid}{m}{
  \ifmmode
    \textsf{\detokenize{#1}}
  \else
    \group_begin:
    \urlstyle{sf}
    \tl_set:Nn \l_tmpa_tl {#1}
    \tl_replace_all:Nnn \l_tmpa_tl {\_} {_}
    \exp_args:NV \path \l_tmpa_tl
    \group_end:
  \fi
}
\ExplSyntaxOff
\providecommand{\tr}{\operatorname{tr}}
\newcommand{\tnleanIssueBase}{https://github.com/LionSR/TNLean/issues/}
\newcommand{\tnleanPrBase}{https://github.com/LionSR/TNLean/pull/}
\newcommand{\tnleanDocBase}{https://sirui-lu.com/TNLean/docs/}
\newcommand{\ghissue}[1]{\href{\tnleanIssueBase#1}{GitHub issue \##1}}
\newcommand{\ghpr}[1]{\href{\tnleanPrBase#1}{PR \##1}}
\newcommand{\doclink}[1]{\href{\tnleanDocBase#1.html}{\path{#1}}}

% Dirac notation and the identity operator, as in blueprint/src/macros/common.tex.
% \providecommand keeps note-local definitions authoritative.
\usepackage{dsfont}
\providecommand{\ket}[1]{|#1\rangle}
\providecommand{\bra}[1]{\langle #1|}
\providecommand{\braket}[2]{\langle #1 | #2 \rangle}
\providecommand{\ketbra}[2]{|#1\rangle\!\langle #2|}
\providecommand{\Id}{\mathds{1}}
\providecommand{\one}{\mathds{1}}
\providecommand{\C}{\mathbb{C}}
\providecommand{\MN}[1]{M_{#1}(\C)}
\providecommand{\MD}{M_D(\C)}

% External referencing for paper sources.  The build helper writes sanitized
% auxiliary files under build/paper-aux/ and excludes duplicated source labels.
\usepackage{xr}

% Import a generated paper auxiliary file with a caller-supplied label prefix.
% The candidate paths support builds from docs/paper-gaps/, the repository root,
% and build/paper-aux/ itself.
\newcommand{\importpaperaux}[2]{%
  \IfFileExists{../../build/paper-aux/#2.aux}{%
    \externaldocument[#1]{../../build/paper-aux/#2}%
  }{%
    \IfFileExists{build/paper-aux/#2.aux}{%
      \externaldocument[#1]{build/paper-aux/#2}%
    }{%
      \IfFileExists{#2.aux}{%
        \externaldocument[#1]{#2}%
      }{}%
    }%
  }%
}

% General external reference macros
\newcommand{\paperref}[2]{\ref{#1-#2}}
\newcommand{\papereqref}[2]{(\ref{#1-#2})}

% CPSV16 source (1606.00608 / MPDO-22-12-17-2.tex) external document and shortcuts
\importpaperaux{cpsv16-}{1606.00608/MPDO-22-12-17-2-tnlean}
\newcommand{\cpsvref}[1]{\paperref{cpsv16}{#1}}
\newcommand{\cpsveqref}[1]{\papereqref{cpsv16}{#1}}

% Verdict marker: \gapnote{<kind>}{<status>}. Kind is the policy
% classification (clarification, local-correction, scope-restriction,
% unfaithful, false-source, open-gap); status is open, wip (elimination
% underway), resolved, or historical. Machine-read by the site index;
% prints a verdict line.
\newcommand{\gapnote}[2]{\par\noindent\textbf{Verdict:} #1 (#2).\par}


\title{Paper Comparison for BNT Inputs in the Non-periodic MPS Fundamental Theorem}
\author{}
\date{2026-05-06}

\begin{document}
\maketitle
\gapnote{open-gap}{open}

\section{Source theorem and comparison standard}

The nonperiodic matrix product vector Fundamental Theorem concerns tensors in
canonical form
\cite{Cirac2016MPDO_arXiv,Cirac2017MPDO_AnnPhys,Cirac2021Matrix}. It is not a
theorem about arbitrary block-diagonal tensors supplemented by unrelated
finite-length assumptions. The finite-length hypotheses isolated in the
formalization name ingredients of the source proof: canonical-form reduction,
selection of a basis of normal tensors (BNT), fixed-length block-injective span
after blocking, phase-gauge matching, and comparison of power sums.

This note uses the cited sources, rather than issue numbers, as the comparison
standard. Issue trackers may schedule the formal work, but they do not determine
the mathematical boundary. A conditional formal theorem matches the source
only if each additional hypothesis is one of the source ingredients identified
here, or if the theorem explicitly records that the corresponding source
ingredient has not yet been proved.

For a tensor $A=(A^i)_i$, define its matrix product vector at positive length
$N$ by the amplitudes
\begin{align}
    \bra{i_1,\ldots,i_N}V^{(N)}(A)
        &=
        \tr\left(A^{i_1}\cdots A^{i_N}\right),
        \qquad N\geq 1.
    \label{eq:rmp_nonperiodic_bnt_comparison_inputs_positive_length_mpv}
\end{align}
The source uses only positive lengths
\cite[Definition~\texttt{MPV}, lines~130--151]{Cirac2016MPDO_arXiv}. A formal
theorem should therefore not use the empty word as a source hypothesis or
conclusion.

Canonical-form reduction gives
\begin{align}
    A^i
        &=
        \bigoplus_k \mu_k A_k^i,
    \label{eq:rmp_nonperiodic_bnt_comparison_inputs_canonical_blocks}\\
    \sum_k D_k
        &\leq D,
    \label{eq:rmp_nonperiodic_bnt_comparison_inputs_canonical_dimensions}
\end{align}
where $D$ is the original bond dimension and $D_k$ is the dimension of the
$k$th nonzero invariant block
\cite[lines~214--225]{Cirac2016MPDO_arXiv}. Periods are then removed by
blocking. Zero blocks and period-removal data must remain separate from the
later injectivity blocking.

After selecting BNT representatives $A_1,\ldots,A_g$, the source writes
\begin{align}
    A^i
        &=
        X\left[\bigoplus_{j=1}^{g}\left(M_j\otimes A_j^i\right)\right]X^{-1},
    \label{eq:rmp_nonperiodic_bnt_comparison_inputs_bnt_tensor_decomposition}\\
    \ket{V^{(N)}(A)}
        &=
        \sum_{j=1}^{g}
        \left(\sum_{q=1}^{r_j}\mu_{j,q}^{N}\right)
        \ket{V^{(N)}(A_j)}.
    \label{eq:rmp_nonperiodic_bnt_comparison_inputs_bnt_mpv_decomposition}
\end{align}
The comparison concerns these selected representatives, not the raw
cyclic-sector or repeated-copy lists produced during canonical-form reduction.

After the block-injectivity step, the blocked evaluations $\widetilde A^i$
satisfy
\begin{align}
    \operatorname{span}\{\widetilde A^i\}_i
        &=
        \bigoplus_{j=1}^{g} M_{D_j}(\C).
    \label{eq:rmp_nonperiodic_bnt_comparison_inputs_block_injective_span}
\end{align}
The formal proof must establish the homogeneous fixed-length direct-sum span in
(\ref{eq:rmp_nonperiodic_bnt_comparison_inputs_block_injective_span}), or retain
it explicitly as a hypothesis.

Once the BNT representatives of two tensors have been matched, the equal case
compares their weights through
\begin{align}
    \sum_q \mu_{j,q}^{N}
        &=
        \sum_q \left(\nu_{j,q}e^{\mathrm{i}\phi_j}\right)^{N},
        \qquad\text{for every matched $j$ and every $N\geq 1$}.
    \label{eq:rmp_nonperiodic_bnt_comparison_inputs_equal_case_power_sums}
\end{align}
The finite power-sum lemma is applied only after this representative-level
matching.

\section{Source labels and their formal meaning}

The local 2016 source labels the canonical-form reduction
\texttt{eq:II\_Aiplusk1}:
\begin{align}
    A^i
        &=
        \sum_{k=1}^{r}P_k B^iP_k
    \label{eq:rmp_nonperiodic_bnt_comparison_inputs_source_invariant_reduction}\\
        &=
        \bigoplus_{k=1}^{r}\mu_k A_k^i.
    \label{eq:rmp_nonperiodic_bnt_comparison_inputs_source_canonical_reduction}
\end{align}
The sentence following that source display states
$\sum_k D_k\leq D$, allowing an all-zero leftover block
\cite[lines~214--219]{Cirac2016MPDO_arXiv}. The source display
\texttt{eq:II\_CF1} restates the canonical tensor as
\begin{align}
    A^i
        &=
        \bigoplus_{k=1}^{r}\mu_k A_k^i.
    \label{eq:rmp_nonperiodic_bnt_comparison_inputs_source_canonical_form}
\end{align}

The phrase ``zero tail'' is formalization terminology. It denotes the unused
all-zero summand of dimension
\begin{align}
    D_{\mathrm{zero}}
        &=
        D-\sum_k D_k.
    \label{eq:rmp_nonperiodic_bnt_comparison_inputs_zero_tail_dimension}
\end{align}
It is not terminology from the source and should not appear without this
definition.

The source phase-gauge relation \texttt{eq:II:A=XAX} is
\begin{align}
    \widetilde A_k
        &=
        e^{\mathrm{i}\phi_k}X_kA_jX_k^{-1}.
    \label{eq:rmp_nonperiodic_bnt_comparison_inputs_source_phase_gauge}
\end{align}
Together with the BNT characterization \texttt{prop:char-BNT}, this relation
identifies normal blocks modulo phase and invertible gauge
\cite[lines~264--280 and 1135--1148]{Cirac2016MPDO_arXiv}. The source BNT is
therefore a list of representatives of phase-gauge equivalence classes. A
common phase cover on raw cyclic sectors is not yet the source matching
statement; it must first descend to the selected representatives.

The source subequation \texttt{eq:II\_ABasicTensors} expands the BNT data as
\begin{align}
    A^i
        &=
        X\left[\bigoplus_{j=1}^{g}\left(M_j\otimes A_j^i\right)\right]X^{-1}
    \label{eq:rmp_nonperiodic_bnt_comparison_inputs_source_bnt_first_form}\\
        &=
        \bigoplus_{j=1}^{g}\bigoplus_{q=1}^{r_j}
        \mu_{j,q}X_{j,q}A_j^iX_{j,q}^{-1}.
    \label{eq:rmp_nonperiodic_bnt_comparison_inputs_source_bnt_copy_form}
\end{align}
Its matrix product vector expansion is the source display
\texttt{decBSV}. The representative index $j$ and copy index $q$ have
different roles: repeated blocks contribute through the power sums attached to
one representative, as in
(\ref{eq:rmp_nonperiodic_bnt_comparison_inputs_bnt_mpv_decomposition}). They
should not become theorem-level sector-matching data unless the theorem
specifically recovers their multiplicities or weights.

The source labels the block-injective definition and its bound
\texttt{defnbi} and \texttt{propblockinj}. These statements give an exact
physical-length span after blocking, namely
(\ref{eq:rmp_nonperiodic_bnt_comparison_inputs_block_injective_span}).
Cumulative algebra fullness and length-zero membership of the identity do not
give this conclusion.

The proportional Fundamental Theorem is \texttt{thm1}, and the equal-MPV
corollary is \texttt{II\_cor2}
\cite{Cirac2016MPDO_arXiv}. In the proof appendix,
\texttt{equalMPS} gives the normal-block overlap and phase-gauge criterion,
\texttt{eqV} gives the phase-equality corollary, and
\texttt{Lem:app\_simple} is the finite power-sum lemma
\cite[lines~1155--1163]{Cirac2016MPDO_arXiv}. These are the source objects that
the formal BNT comparison must reproduce. Raw flattened cyclic sectors, common
phase covers, and cumulative pair algebras are formal intermediates. They
support the source theorem only after they are proved to yield the labelled BNT
objects.

These observations impose five translation rules.

\begin{enumerate}
    \item The source displays \texttt{eq:II\_Aiplusk1} and
    \texttt{eq:II\_CF1} give nonzero blocks whose dimensions sum to at most the
    original bond dimension. The formal zero tail is only the omitted all-zero
    summand. It is not a normal tensor, a BNT representative, or an additional
    scalar weight.

    \item The source relation \texttt{eq:II:A=XAX} and characterization
    \texttt{prop:char-BNT} quotient canonical blocks by phase and gauge
    \cite[lines~264--280 and 1135--1148]{Cirac2016MPDO_arXiv}. A matching
    statement for raw cyclic sectors must descend to these equivalence classes
    before it matches the source BNT theorem.

    \item The source formulas \texttt{eq:II\_ABasicTensors} and
    \texttt{decBSV} separate representatives from their repeated copies.
    Repetitions enter through coefficient power sums, not through an additional
    family of unmatched source sectors.

    \item The source statements \texttt{defnbi} and
    \texttt{propblockinj} require exact-length direct-sum span. They do not
    follow from cumulative fullness or from the empty word.

    \item The source theorem \texttt{thm1} first matches BNT representatives.
    The equal-case corollary \texttt{II\_cor2} then applies
    \texttt{Lem:app\_simple} to their coefficient power sums
    \cite{Cirac2016MPDO_arXiv}. The formal equal-case proof must preserve this
    order.
\end{enumerate}

These rules also determine theorem naming in this part of the formalization.
Results that only package raw cyclic-sector data into source hypotheses are
supporting lemmas. Theorems should describe the source-level milestones or
state the remaining source hypotheses explicitly.

\section{Canonical form, period removal, and BNT selection}

The source starts with a tensor $B$, removes invariant subspaces, and replaces
it by the block-diagonal tensor in
(\ref{eq:rmp_nonperiodic_bnt_comparison_inputs_canonical_blocks}), which gives
the same positive-length matrix product vectors
\cite[lines~214--225]{Cirac2016MPDO_arXiv}. It then studies the completely
positive map of each block. Blocking a common multiple of the periods removes
periodic peripheral eigenvalues
\cite[lines~227--231]{Cirac2016MPDO_arXiv}. The resulting blocks are normal in
the source sense, and the tensor is in canonical form
\cite[lines~233--255]{Cirac2016MPDO_arXiv}. The review gives the same sequence:
block diagonalization, period-removing blocking, primitive transfer maps, and
canonical form \cite[lines~1793--1839]{Cirac2021Matrix}.

Period removal and finite-length injectivity are distinct operations. The first
makes the peripheral spectrum primitive. The second supplies injective or
block-injective span after normality is available.

For a canonical tensor, a BNT is a family of normal tensors
$A_1,\ldots,A_g$ whose matrix product vectors span those of $A$ at every
length and are eventually linearly independent
\cite[lines~271--274]{Cirac2016MPDO_arXiv}. Every canonical normal block is a
phase times an invertible conjugate of one selected tensor, and no two selected
tensors are phase-gauge equivalent
\cite[lines~278--280]{Cirac2016MPDO_arXiv}. The review uses the same definition
and characterization \cite[lines~1846--1859]{Cirac2021Matrix}.

The BNT decomposition is
(\ref{eq:rmp_nonperiodic_bnt_comparison_inputs_bnt_mpv_decomposition})
\cite[lines~283--302]{Cirac2016MPDO_arXiv};
\cite[lines~1864--1885]{Cirac2021Matrix}. Hence the source compares selected
representatives. Repeated copies and cyclic representatives contribute through
the power sums of their weights.

For two canonical tensors $A$ and $B$, the source Fundamental Theorem states
that proportional matrix product vector families have BNTs related by a
permutation, phases, and invertible gauge matrices
\cite[lines~347--352]{Cirac2016MPDO_arXiv};
\cite[lines~1890--1900]{Cirac2021Matrix}. In the equal case, the source obtains
(\ref{eq:rmp_nonperiodic_bnt_comparison_inputs_equal_case_power_sums}) for each
matched representative. The finite power-sum lemma
\texttt{Lem:app\_simple} recovers the multiplicities and weights
\cite[lines~1155--1163]{Cirac2016MPDO_arXiv}. The proof of
\texttt{II\_cor2} then assembles the blockwise gauges into one global
conjugating matrix \cite[lines~1184--1192]{Cirac2016MPDO_arXiv}.

\section{Fixed-length span after blocking}

The source uses two blocking operations. Period-removing blocking produces
primitive normal blocks. A later quantum Wielandt step produces finite-length
injectivity \cite{Sanz2010Quantum}. These operations should not be merged.

For one normal tensor of bond dimension $D$, the 2016 source states that at
most $D^4$ blockings make the tensor injective
\cite[line~332]{Cirac2016MPDO_arXiv}. The earlier representation paper defines
the length-$L$ map
\begin{align}
    \Gamma_L(X)
        &=
        \sum_{i_1,\ldots,i_L}
        \tr\left(XA_{i_1}\cdots A_{i_L}\right)
        \ket{i_1,\ldots,i_L}.
    \label{eq:rmp_nonperiodic_bnt_comparison_inputs_gamma_length_map}
\end{align}
Injectivity of $\Gamma_L$ is equivalent to the length-$L$ word products
spanning the entire matrix algebra
\cite[lines~888--908]{PerezGarcia2007MPSReps}. This is a homogeneous
fixed-length statement at a positive blocking length, not membership in a
cumulative span. The review records a sharper MPS-specific bound
\cite[lines~1942--1945]{Cirac2021Matrix}; the comparison here concerns the form
of the conclusion, not the optimal numerical bound.

For several BNT blocks, block injectivity means that one physical tensor spans
all matrix blocks simultaneously:
\begin{align}
    \operatorname{span}
        \left\{\left(A_1^i,\ldots,A_g^i\right)\right\}_i
        &=
        \bigoplus_{j=1}^{g}M_{D_j}(\C)
        \qquad\text{after a common blocking}.
    \label{eq:rmp_nonperiodic_bnt_comparison_inputs_common_block_span}
\end{align}
The 2016 source cites the P\'erez-Garc\'ia--Verstraete--Wolf--Cirac bound: if
each normal tensor is injective at length $L$, then
$3(L+1)(D-1)$ blockings make the canonical tensor block-injective
\cite[lines~317--345]{Cirac2016MPDO_arXiv,PerezGarcia2007MPSReps}. Combining
this with $L\leq D^4$ gives the stated $3D^5$ bound. The review records the
same conclusion and bound \cite[lines~2114--2124]{Cirac2021Matrix}.

The earlier representation paper supplies the direct-sum input
\cite{PerezGarcia2007MPSReps}. Its condition $C1$ is exact-length
injectivity of $\Gamma_{L_0}$, equivalently the span of all length-$L_0$
word products
\cite[lines~888--908]{PerezGarcia2007MPSReps}. It fixes canonical blocks
satisfying $C1$, defines
\begin{align}
    L_0
        &=
        \max_j L_0^{A^j},
    \label{eq:rmp_nonperiodic_bnt_comparison_inputs_common_injectivity_length}
\end{align}
orders their bond dimensions, and assumes that the corresponding block states
are pairwise different
\cite[lines~1321--1330]{PerezGarcia2007MPSReps}. A matrix factorization lemma
and direct-sum lemma then make the sum of the length-$L$ MPS subspaces direct
whenever
\begin{align}
    L
        &\geq
        3(b-1)(L_0+1).
    \label{eq:rmp_nonperiodic_bnt_comparison_inputs_direct_sum_length_bound}
\end{align}
This argument appears at
\cite[lines~1333--1408]{PerezGarcia2007MPSReps}. The source input is therefore
a homogeneous fixed-length direct-sum span, or its trace-dual separation form.
The note
\path{docs/paper-gaps/pgvwc07_direct_sum_input.tex} records it separately
because it is a cited proof component rather than a named theorem of the 2016
source.

Block injectivity is stronger than separate injectivity of the blocks. It
provides the cross-block separation needed to prescribe distinct matrices in
different BNT components at one common length. In trace-dual form, for one
fixed positive length $L$,
\begin{align}
    &\left[\forall w,\ |w|=L,\quad
        \sum_{j=1}^{g}\tr\left(\Delta_jA_j^w\right)=0\right]
    \notag\\
    &\qquad\Longrightarrow
        \Delta_j=0
        \quad\text{for every $j$}.
    \label{eq:rmp_nonperiodic_bnt_comparison_inputs_trace_dual_separation}
\end{align}
Equivalently,
\begin{align}
    \operatorname{span}
        \left\{\left(A_1^w,\ldots,A_g^w\right): |w|=L\right\}
        &=
        \bigoplus_{j=1}^{g}M_{D_j}(\C).
    \label{eq:rmp_nonperiodic_bnt_comparison_inputs_exact_length_direct_sum}
\end{align}
The pair trace-separation predicates in the nonperiodic proof express this dual
form. They are not separate paper theorems. In the later MPDO argument, the
source defines
\begin{align}
    C_L(V)
        &=
        \operatorname{span}(V^L)
    \label{eq:rmp_nonperiodic_bnt_comparison_inputs_source_exact_length_space}\\
        &=
        \mathcal A_1
        \qquad\text{for a suitable fixed $L$}.
    \label{eq:rmp_nonperiodic_bnt_comparison_inputs_source_exact_length_algebra}
\end{align}
The block-injective proposition supplies the second equality
\cite[lines~1984--1988]{Cirac2016MPDO_arXiv}.

\section{Homogeneous padding}

The injectivity and block-injectivity conclusions are homogeneous statements at
a specified positive length. They do not follow from the empty word belonging
to a cumulative word span.

For a one-letter alphabet in bond dimension one, let
\begin{align}
    A_0
        &=
        0,
    \label{eq:rmp_nonperiodic_bnt_comparison_inputs_zero_pair_first_tensor}\\
    B_0
        &=
        0.
    \label{eq:rmp_nonperiodic_bnt_comparison_inputs_zero_pair_second_tensor}
\end{align}
The length-one pair span is generated by the zero pair, whereas the identity
pair occurs at length zero. Thus cumulative identity membership does not imply
identity membership in later homogeneous spans.

Cumulative fullness is also insufficient. In bond dimension one, consider the
one-letter pair word values
\begin{align}
    \left(A^n,B^n\right)
        &=
        \left(2^n,3^n\right).
    \label{eq:rmp_nonperiodic_bnt_comparison_inputs_scalar_pair_words}
\end{align}
The cumulative span through length one is all of $\C\oplus\C$, because
$(1,1)$ and $(2,3)$ are linearly independent. At every positive
homogeneous length $n$, however, the span is only
\begin{align}
    \operatorname{span}
        \left\{\left(2^n,3^n\right)\right\},
    \label{eq:rmp_nonperiodic_bnt_comparison_inputs_scalar_homogeneous_line}
\end{align}
which does not contain $(1,1)$. The source hypotheses therefore cannot be
weakened to cumulative pair-span fullness.

The required conclusion is fixed-length block-injective direct-sum span after
the prescribed blocking. A homogeneous-padding argument based on irreducibility
could prove an equivalent exact-length separation statement, but it would have
to retain the BNT selection, non-phase-gauge-equivalence, and normalization
data. Full cumulative pair span and empty-word membership alone do not suffice.

\section{Comparison after common blocking}

The equal-case source proof begins after canonical form and BNT selection. The
formal after-blocking route must additionally discard the all-zero tail, choose
one common positive physical blocking length, identify nested blocked words
with flat words, distinguish period removal from injectivity blocking, and
compare the resulting nonzero sector families.

Suppose a block has period-removal length $m_k$, and let the final common
length be
\begin{align}
    p
        &=
        m_ke_k.
    \label{eq:rmp_nonperiodic_bnt_comparison_inputs_common_blocking_factorization}
\end{align}
Blocking first by $m_k$ and then by $e_k$ gives nested physical labels,
whereas blocking once by $p$ gives flat word labels. Source notation
identifies these alphabets silently. The formal theorem must prove the required
reindexing.

The BNT comparison data must be constructed for the final nonzero sector
families after all blockings and reindexings. The proof must group common
primitive sectors into representative BNT classes and account for repeated
cyclic representatives and transported weights through the power-sum
comparison. Common phase covers and proportional decompositions match the
source only after this grouping. A raw flattened cyclic-sector list may contain
several cyclic representatives of one original block with the same transported
weight, so it is not the BNT used by the source theorem.

\section{Formal translation}

The source comparison gives the following formal requirements.

\begin{itemize}
    \item The formal comparison must distinguish positive-length matrix product
    vector equality from any auxiliary length-zero convention. If a proof uses
    the empty word, it must remove that convention before invoking the source
    theorem.

    \item Canonical-form reduction has two blocking roles. Period-removal
    blocking produces primitive normal blocks; later quantum Wielandt or
    block-injectivity blocking produces exact finite-length span
    \cite{Sanz2010Quantum}.

    \item BNT trace separation is the dual of the fixed-length direct-sum span
    supplied by block-injective canonical form. The formal proof must establish
    this span for selected BNT representatives or retain it explicitly as a
    hypothesis.

    \item Phase-cover and block-matching data must follow from representative
    BNT matching and the equal-case power-sum comparison. They must not enter
    as unrelated assumptions on raw flattened cyclic sectors.

    \item The after-blocking proof needs positive dimensions, injectivity,
    normalization, self-overlap limits, off-overlap decay for inequivalent
    representatives, and equality of finite-length matrix product vector spans.
    Auxiliary structures that package these facts are formal devices; the
    mathematical obligation is to prove them for the representatives selected
    by canonical-form reduction and BNT construction.
\end{itemize}

\section{Current formal anchors}

The one-copy normal-canonical BNT surface is represented by
\leanid{MPSTensor.IsNormalCanonicalFormBNT}, together with explicit separated-block
hypotheses used by the BNT comparison lemmas.\footnote{%
\path{TNLean/MPS/BNT/Construction.lean}.}
This already-grouped formulation is restricted: it suppresses the repeated-copy
data $r_j$ and $\mu_{j,q}$ appearing in the source BNT expansion and appendix
characterization \cite{Cirac2016MPDO_arXiv}.

The former coefficient-array structures and conditional proportional comparison
theorems have been deleted. Their one-shot bounded coefficient assumptions did
not encode the source residual peeling or power-sum comparison. The former
conditional after-blocking matching layer has also been removed. It had
collected comparison data for common primitive nonzero-sector families and
concluded BNT sector matching after one positive common blocking length.

The named blocked-normal-form records and common-primitive block-matching
boundary have likewise been removed. The live anchors are the common-sector
transport construction and the BNT sector-data comparison theorem.\footnote{%
\path{TNLean/MPS/CanonicalForm/SectorComparison/CommonSectorTransport.lean};
\path{TNLean/MPS/FundamentalTheorem/SectorBNT/Fundamental.lean}.}

\section{Named formal gaps}

The nonperiodic proof route retains three obligations from the source argument.

\begin{enumerate}
    \item Construct common representative-level BNT cover data for the
    collapsed nonzero-sector families selected by canonical-form reduction and
    the BNT characterization. This is the representative version of the source
    phase-gauge matching theorem, not a condition on every raw block in an
    intermediate cyclic decomposition
    \cite{Cirac2016MPDO_arXiv,Cirac2021Matrix}.

    \item Prove homogeneous trace separation for distinct ordered pairs of BNT
    blocks. In source notation, this is the trace-dual form of the
    block-injective span
    \begin{align}
        \operatorname{span}
            \left\{\left(A_1^w,\ldots,A_g^w\right): |w|=L\right\}
            &=
            \bigoplus_j M_{D_j}(\C)
    \label{eq:rmp_nonperiodic_bnt_comparison_inputs_named_gap_direct_sum}
    \end{align}
    at one common positive length $L$. It is stronger than fullness of the
    cumulative algebra generated by a pair representation.

    \item Transport this fixed-length separation through the after-blocking
    comparison. Period-removal blocking, injectivity blocking, and the
    nested-to-flat word reindexing must all be accounted for before the
    BNT-level Fundamental Theorem is applied
    \cite{Cirac2016MPDO_arXiv,Cirac2021Matrix}.
\end{enumerate}

\section{Source files checked}

The comparison uses the repository's \LaTeX{} sources, not only published
theorem names.

\begin{itemize}
    \item \path{Papers/1606.00608/MPDO-22-12-17-2.tex}: lines~130--151 define
    positive-length matrix product vectors; lines~214--255 give canonical form
    and period removal; lines~271--302 define and use BNTs; lines~317--345
    define injectivity and block injectivity and give the $3D^5$ bound;
    lines~347--360 state the Fundamental Theorem and equal-case corollary;
    lines~1155--1163 state \texttt{Lem:app\_simple}; lines~1184--1192 compare
    the power sums and assemble the global gauge; lines~1984--1988 use
    $C_L(V)=\operatorname{span}(V^L)$ as an exact-length span
    \cite{Cirac2016MPDO_arXiv}.

    \item \path{Papers/quant-ph_0608197/MPSarchive.tex}: lines~742--763 state
    the translation-invariant canonical form; lines~888--908 define
    $\Gamma_L$ and identify injectivity with the span of length-$L$ word
    products \cite{PerezGarcia2007MPSReps}.

    \item \path{Papers/1210.6613/UncleHMPS.tex}: lines~247--278 record the
    standard form after blocking, including the block-diagonal span condition
    and its injective and block-injective special cases.

    \item \path{Papers/2011.12127/TN-Review-main.tex}: lines~1793--1900 review
    canonical form, BNTs, and the Fundamental Theorem; lines~1942--1945 give
    the normal-tensor Wielandt bound; lines~2114--2124 state the
    block-injective canonical-form bound \cite{Cirac2021Matrix}.
\end{itemize}

\section{Present formal boundary}

The comparison with the 2016 source separates the mathematical layers of its
proof from the history of the implementation. The formalization has established
the exact-length span input for an already-injective BNT surface. It has also
formalized the quantum Wielandt blocking input needed to make normal blocks
injective after further blocking. The remaining work is to transport the final
reblocked BNT cover and formalize the coefficient comparison performed through
power sums in the source \cite{Cirac2016MPDO_arXiv}.

For an already-injective canonical BNT family, identity padding and pairwise
trace separation are no longer separate paper assumptions. The formal result
specializes the three-block direct-sum argument to selected BNT blocks that are
one-site injective. One-site injectivity supplies the fixed length-two and
length-six span inputs used by the selector construction. This theorem proves
the required exact-length direct-sum span on that restricted surface; it does
not prove the general result for an arbitrary pre-existing injectivity length
$L_0$. Passing from normal canonical blocks to this injective surface remains
a separate quantum Wielandt blocking step \cite{Sanz2010Quantum}.

That blocking step is now formalized in the canonical-form setting. A
left-canonical normal tensor has a positive injective blocking length bounded by
$D^4$, and a finite normal-CF-BNT family has one common positive blocking
length for all blocks. These results remove injectivity as a mathematical gap.
They do not identify the iterated blocked alphabet with the direct alphabet used
by the after-blocking comparison. The remaining transport must carry BNT
separation and zero-tail bookkeeping through the refinement of the blocking
length.

The former proportional-decomposition coefficient-limit surface was deleted
rather than treated as a source theorem. The source compares BNT weights by
power sums \cite{Cirac2016MPDO_arXiv}. This remains a genuine proof obligation:
the block-diagonal expansion contains the length-dependent coefficients
\begin{align}
    c_j(N)
        &=
        \mu_j^N,
    \label{eq:rmp_nonperiodic_bnt_comparison_inputs_weight_length_power}
\end{align}
and nonzero weights alone do not imply nonzero full-sequence limits. A faithful
proof must formalize the relevant power-sum and phase comparison before
constructing the \leanid{MPSTensor.BlockPermutationGaugePhaseConclusion} data and
applying the BNT sector-data comparison.

Zero-tail equality and common-sector relabelling are primarily formal
bookkeeping. The source uses positive lengths and treats zero blocks as unused
bond dimension \cite{Cirac2016MPDO_arXiv}. The $N=0$ convention in
\leanid{MPSTensor.SameMPV₂} forces an explicit zero-tail identity, and the
common-sector construction must track casts between grouped and flat blocked
words. These are not additional paper hypotheses, but the final after-blocking
theorem must transport the identities through the same reblocking coordinates
as the nonzero BNT sectors.

The remaining main boundary is therefore precise. Starting from normal-CF-BNT
sectors produced by canonical-form reduction, the proof must transport blocked
non-phase-gauge equivalence, zero-tail identities, and the common injective
blocking length through the nested-to-flat reblocking map. It must then apply
the BNT sector-data comparison.\footnote{The live target is
\leanid{MPSTensor.ft_sector_bnt_equal_sector_data}. The available inputs are
the bounded and common injective-blocking theorems for normal-CF-BNT blocks and
the already-injective direct-sum trace-separation theorems.}
The source suppresses this bookkeeping in its notation; the formal proof must
make it explicit.

\section{Verdict and closing criterion}

The conditional formal hypotheses do not create a mathematical discrepancy
when they name source inputs still to be derived from canonical form. The
source assumes canonical tensors and compares their BNTs; it does not assume
arbitrary phase covers, raw cyclic-sector matching, cumulative pair-span
fullness, or length-zero matrix product vector equality
\cite{Cirac2016MPDO_arXiv,Cirac2021Matrix}.

The present boundary is therefore a proof gap rather than a different theorem.
The remaining inputs are the representative-level BNT cover, transport of the
exact-length direct-sum span through common blocking and word reindexing, and
the coefficient comparison by power sums. Until these are proved, the
downstream after-blocking theorem should remain visibly conditional.

A proof obligation is discharged only when three statements agree. The source
assertion must be identified as canonical form, BNT selection,
block-injective fixed-length span, phase-gauge matching, or power-sum
comparison. The formal theorem must state that assertion for the same blocked
and reindexed families, or retain the unresolved hypothesis explicitly. The
final comparison theorem must use that assertion rather than an unrelated
assumption of similar shape.

Conditional theorems usefully expose missing inputs. They do not close the
source theorem until those inputs have been proved or the final result has been
deliberately restricted.

\bibliographystyle{alpha}
\bibliography{references}

\end{document}
